\documentclass{beamer}
\usepackage[T1]{fontenc}
\usepackage[utf8]{inputenc}
\usepackage{textcomp}
\usepackage{eurosym}
\DeclareUnicodeCharacter{00B0}{\textdegree}
\DeclareUnicodeCharacter{20AC}{\euro}
\DeclareUnicodeCharacter{2013}{--}
\DeclareUnicodeCharacter{2014}{---}
\DeclareUnicodeCharacter{2018}{`}
\DeclareUnicodeCharacter{2019}{'}
\DeclareUnicodeCharacter{201C}{``}
\DeclareUnicodeCharacter{201D}{''}
\DeclareUnicodeCharacter{2026}{\ldots{}}
\usepackage[spanish,es-noshorthands]{babel}

% Definición del color corporativo aproximado para Pantone Warm Red C
\definecolor{UMWarmRed}{HTML}{F9423A}

% Tema y esquema de colores
\usetheme{Madrid}
\usecolortheme{rose} % Usamos un esquema de color que pueda combinarse bien

% Configuración de colores usando el color corporativo
\setbeamercolor{title}{bg=UMWarmRed, fg=white}
\setbeamercolor{frametitle}{bg=UMWarmRed, fg=white}
\setbeamercolor{structure}{fg=UMWarmRed}
\setbeamercolor{section number projected}{bg=UMWarmRed,fg=white}
\setbeamercolor{itemize item}{fg=UMWarmRed}
\setbeamercolor{itemize subitem}{fg=UMWarmRed}
\setbeamercolor{enumerate item}{fg=UMWarmRed}

\title[DISEÑO DE EXPERIMENTOS]{METODOLOGÍA Y DISEÑO DE EXPERIMENTOS: \\ PRINCIPIOS BÁSICOS}
\author[Alfonso Rosa García]{Alfonso Rosa García }
\institute[G9]{\textbf{Grupo de Universidades G-9}}

\date{9 y 11 de marzo de 2026 \\ 10:00--13:00}

\begin{document}

\begin{frame}
  \titlepage
\end{frame}

% Diapositiva de presentación del profesor
\begin{frame}{Presentación del Profesor}
      \begin{itemize}
         \item Alfonso Rosa García
         \item Profesor de Fundamentos del Análisis Económico.
         \item \textbf{Especialización:} Economía y Finanzas Experimentales, Economía de redes y Teoría de Juegos.
         \item \textbf{Email:} \texttt{alfonso.rosa@um.es}
      \end{itemize}
\end{frame}

% Diapositiva de dinámica del curso
\begin{frame}{Dinámica de las sesiones}
   \begin{block}{Modalidad y Estructura}
      \begin{itemize}
         \item \textbf{Sesiones:} Lunes 9 y miércoles 11 de marzo de 2026, de 10:00 a 13:00.
         \item \textbf{Modalidad:} Presencial/online según la sede del programa G9.
         \item \textbf{Enfoque:} Multidisciplinar, con conceptos generales aplicables a todas las áreas de conocimiento.
         \item \textbf{Evaluación:} Tareas en grupos durante la sesión y test individual a realizar en el plazo de una semana.
         \item \textbf{Asistencia:} Si no puedes asistir o surge algún problema, contacta conmigo a través del campus virtual de vuestra universidad.
      \end{itemize}
   \end{block}
\end{frame}


\begin{frame}{Objetivos del Curso}
\textbf{Objetivos generales}
\begin{itemize}
    \item Reflexionar sobre el proceso de generación de ideas y planificación de la investigación como fases previas al diseño experimental.
    \item Introducir los conceptos básicos de diseño y metodología experimental.
    \item Diferenciar correlación y causalidad con ejemplos prácticos.
    \item Conocer los principios de los diseños experimentales y su aplicación en distintas disciplinas.
    \item Destacar principios estadísticos clave y herramientas para el análisis de datos experimentales.
    \item Fomentar prácticas éticas en la investigación experimental.
\end{itemize}
\end{frame}

\begin{frame}{Programa del Curso}
   \begin{block}{Contenidos (según la guía docente)}
      \begin{enumerate}
         \item Generación de ideas de investigación.
         \item Planificación de la investigación.
         \item Correlación y causalidad.
         \item Tipos de experimentos.
      \end{enumerate}
   \end{block}
   \vspace{0.3cm}
   \textbf{Enfoque del curso:} Nos centraremos especialmente en los puntos 3 y 4, profundizando en el diseño, análisis estadístico y ética de los experimentos. Los puntos 1 y 2 se abordan de forma introductoria como fases previas que conectan con el diseño experimental.
\end{frame}

\begin{frame}{De la idea al experimento}
   \textbf{1. Generación de ideas de investigación}
   \begin{itemize}
      \item Las ideas de investigación surgen de la observación de fenómenos, la revisión de la literatura, lagunas en el conocimiento existente o necesidades prácticas.
      \item Una buena pregunta de investigación es específica, relevante y abordable con los recursos disponibles.
   \end{itemize}
   \vspace{0.3cm}
   \textbf{2. Planificación de la investigación}
   \begin{itemize}
      \item Implica formular hipótesis claras, identificar las variables clave (independientes, dependientes y de control) y seleccionar la metodología adecuada.
      \item La planificación rigurosa es la base sobre la que se construye un buen diseño experimental, que será el foco de este curso.
   \end{itemize}
\end{frame}

\begin{frame}{Conexión con el diseño experimental}
   \centering
   \textbf{El diseño experimental como herramienta para responder preguntas de investigación}
   \vspace{0.4cm}

   \begin{itemize}
      \item Una vez que tenemos una \textbf{idea de investigación} y hemos \textbf{planificado} cómo abordarla, necesitamos herramientas para establecer relaciones causales.
      \item Distinguir \textbf{correlación} de \textbf{causalidad} es fundamental para elegir el método adecuado.
      \item El \textbf{diseño experimental} nos proporciona el marco para controlar variables y obtener evidencia causal.
      \item En este curso veremos los principios, tipos de diseños, herramientas estadísticas y consideraciones éticas que hacen posible una investigación experimental rigurosa.
   \end{itemize}
\end{frame}

\end{document}










